\documentclass[10pt,a4paper]{article}

\usepackage[margin=0.65in]{geometry}
\usepackage[T1]{fontenc}
\usepackage{lmodern}
\usepackage{enumitem}
\usepackage[hidelinks]{hyperref}
\usepackage{titlesec}
\usepackage{xcolor}
\usepackage{qrcode}

%====================
% Color Theme
%====================
\definecolor{PrimaryColor}{HTML}{1F4E79}
\definecolor{SecondaryColor}{HTML}{4B5563}

\setlength{\parindent}{0pt}
\setlist[itemize]{leftmargin=*, noitemsep, topsep=2pt}

\titleformat{\section}
  {\large\bfseries\color{PrimaryColor}}
  {}
  {0em}
  {}
  [\color{SecondaryColor}\titlerule]

\hypersetup{
    colorlinks=true,
    urlcolor=PrimaryColor,
    linkcolor=PrimaryColor
}

\urlstyle{same}

\begin{document}

%====================
% Header with QR
%====================
\begin{minipage}[t]{0.78\textwidth}
    {\LARGE \textbf{\color{PrimaryColor}Abdelrahman Elsaid Elsawy}}\\
    {\color{SecondaryColor}Robotik- und Fahrzeugsimulations-Ingenieur}\\
    ROS/ROS2 \textbar{} BeamNG.tech \textbar{} NVIDIA Isaac Sim \textbar{} RL-Training \textbar{} ONNX \textbar{} Automatisierte QA\\
    E-Mail: \href{mailto:abdelrahman.elsaid.elsawy@gmail.com}{abdelrahman.elsaid.elsawy@gmail.com}\\
    LinkedIn: \href{https://www.linkedin.com/in/abdelrahman-elsaid-elsawy}{linkedin.com/in/abdelrahman-elsaid-elsawy}\\
    Telefon/WhatsApp: \href{https://wa.me/905365231899}{+90 536 523 1899}\\
    Standort: Istanbul, Türkei
\end{minipage}
\hfill
\begin{minipage}[t]{0.18\textwidth}
    \raggedleft
    \qrcode[height=1.7cm]{https://www.linkedin.com/in/abdelrahman-elsaid-elsawy}\\
    {\scriptsize LinkedIn QR}
\end{minipage}

%====================
\section*{Zusammenfassung}

Robotik- und Softwareingenieur mit über 6 Jahren Erfahrung in ROS/ROS2, Systemintegration, autonomen Systemen und hochpräziser Fahrzeugsimulation. Spezialisiert auf die Anbindung von Simulationsplattformen an Roboter- und Engineering-Toolchains, darunter BeamNG.tech, ROS/ROS2, MATLAB und Simulink. Erfahrung in der Entwicklung von API-Brücken, Echtzeit-Simulationssteuerung, Datenerfassung, automatisiertem QA-Testing, Regressionsanalyse, Fehleruntersuchung, technischem Support und kundenorientiertem Engineering für B2B-Simulations-Workflows. Weitere Schwerpunkte sind NVIDIA Isaac Sim, Reinforcement-Learning-Trainingsworkflows, ONNX-Bereitstellung von Modellen und ANN-basierte Analyseforschung mit MATLAB.

%====================
\section*{Kernkompetenzen}

\textbf{\color{PrimaryColor}Robotik:} ROS, ROS2, SLAM, MoveIt, Gazebo, RViz, autonome Navigation, Sensordatenfusion\\
\textbf{\color{PrimaryColor}Simulation:} BeamNG.tech, NVIDIA Isaac Sim, Fahrzeugsimulation, hochpräzise Simulation, Simulations-Workflows\\
\textbf{\color{PrimaryColor}KI und Modellbereitstellung:} Reinforcement Learning, RL-Trainingsworkflows, ONNX, künstliche neuronale Netze, ANN-basierte Analyse\\
\textbf{\color{PrimaryColor}Programmierung:} C, C++, Python, Lua, SQL, Bash, XML, JSON\\
\textbf{\color{PrimaryColor}Integration:} API-Brücken, Systemintegration, Echtzeit-Simulationssteuerung, Datenerfassung, MATLAB, Simulink\\
\textbf{\color{PrimaryColor}Testing und QA:} automatisierte QA-Tests, Regressionstests, Bug-Reproduktion, Root-Cause-Analyse, Fehleruntersuchung\\
\textbf{\color{PrimaryColor}Hardware:} Arduino, ESP32, Raspberry Pi, LiDAR, IMU, Ultraschall, SPI/I2C/UART, PCB-Design, Löten, Prototyping, Catia\\
\textbf{\color{PrimaryColor}Werkzeuge:} Git, Docker, Ubuntu, Zendesk, API-Dokumentation

%====================
\section*{Berufserfahrung}

\textbf{\color{PrimaryColor}Research Software Engineer} \hfill Nov 2021 -- heute\\
\textbf{BeamNG GmbH} \hfill Remote

\begin{itemize}
    \item Integration von BeamNG.tech mit ROS/ROS2, MATLAB und Simulink für hochpräzise Fahrzeugsimulation, Robotik-Workflows und Entwicklung autonomer Systeme.
    \item Entwicklung und Wartung von API-Brücken für Echtzeit-Simulationssteuerung, Fahrzeuginteraktion, Datenerfassung und externe Tool-Integration.
    \item Durchführung von automatisierten QA-Tests, Regressionsanalysen und Fehleruntersuchungen mit C++, Python und Lua zur Verbesserung von Engine-Stabilität, Simulationszuverlässigkeit und Integrationsqualität.
    \item Untersuchung von kunden- und intern gemeldeten technischen Problemen durch Bug-Reproduktion, Analyse des API-Verhaltens und der Szenariostrukturen, Identifikation von Root Causes sowie Bereitstellung von Fixes, Workarounds oder Implementierungshinweisen.
    \item Unterstützung globaler B2B-Kunden über Zendesk bei Simulation-, API-, ROS/ROS2-, MATLAB/Simulink- und Integrationsproblemen.
    \item Verbesserung der Entwicklerdokumentation für BeamNG.tech-API-Nutzung, MATLAB/Simulink-Workflows, Tutorials, HIL/DIL-Material und Integrationshinweise.
    \item Überführung wiederkehrender QA- und Support-Erkenntnisse in wiederverwendbare Testfälle, Doku-Verbesserungen und umsetzbares Produktfeedback zur Steigerung der Entwicklererfahrung.
    \item Beiträge zu den Engineering-Repositories von BeamNG.tech, einschließlich BeamNGpy, API-Dokumentation, Simulations-Workflows, QA-Fixes, Tutorials und Integrationsbeispiele.
    \item Implementierung und Validierung von Szenario- und Missionslade-Workflows, einschließlich programmatischer Szenario-Erkennung, Szenario-Laden, Fahrzeugabruf, Sensoranbindung, Fahrzeugverbindung, skriptgesteuerter Steuerung und Sensorabfrage.
\end{itemize}

\textbf{\color{PrimaryColor}Mechatronics Systems Engineer} \hfill Feb 2020 -- Jul 2021\\
\textbf{Ideasparq Robotics Sdn Bhd} \hfill Vor Ort

\begin{itemize}
    \item Entwicklung von ROS-basierten Steuerungssystemen und autonomen Navigationslösungen für Serviceroboter.
    \item Implementierung von Sensordatenfusionsalgorithmen und Test von Robotikplattformen mit Gazebo, RViz, Docker und Ubuntu-basierten Robotik-Workflows.
    \item Arbeit an der Robotik-Systemintegration über Software, Sensoren, Steuerungslogik, Simulation und Deployment-Umgebungen hinweg.
\end{itemize}

\textbf{\color{PrimaryColor}Autonomous Robotics Engineer} \hfill Jun 2019 -- Jan 2020\\
\textbf{ENDSTRUCT} \hfill Vor Ort

\begin{itemize}
    \item Entwicklung autonomer Navigations- und Zustandsüberwachungssysteme für mobile GOVTECH-Roboter.
    \item Arbeit an mobiler Robotik-Software in den Bereichen Navigation, Überwachung und feldorientiertes Robotersystemverhalten.
\end{itemize}

\textbf{\color{PrimaryColor}Research Assistant} \hfill Sep 2019 -- Dez 2019\\
\textbf{Autonomous Vehicle Laboratory}

\begin{itemize}
    \item Verarbeitung von 3D-LiDAR-Daten für Echtzeit-Objekterkennung.
    \item Entwicklung einer ROS-basierten Middleware zur Sensorabstraktion und für Forschungsworkflows autonomer Fahrzeuge.
    \item Anwendung künstlicher neuronaler Netze für ingenieurwissenschaftliche Analyseforschung und modellbasierte Vorhersage-Workflows.
\end{itemize}

\textbf{\color{PrimaryColor}Praktikant (Mechatronik)} \hfill Jun 2018 -- Sep 2018\\
\textbf{University Malaysia Pahang}

\begin{itemize}
    \item Praxiserfahrung mit TurtleBot3 und ROS-Robotersimulation gesammelt.
\end{itemize}

%====================
\section*{Ausgewählte Technische Beiträge}

\begin{itemize}
    \item BeamNG.tech-Integration mit ROS/ROS2, MATLAB und Simulink.
    \item Entwicklung von BeamNGpy und API-Brücken für Echtzeit-Simulationssteuerung und Datenzugriff.
    \item Validierung von Szenario- und Missionslade-Workflows für automatisierte Simulations-Workflows.
    \item QA-Automatisierung, Regressionstests, Bug-Reproduktion und Root-Cause-Analyse.
    \item Kundenorientierter Engineering-Support für globale B2B-Simulationsnutzer.
    \item Dokumentationsverbesserungen für APIs, Tutorials, Integrations-Workflows, MATLAB/Simulink und HIL/DIL-Material.
    \item Robotik-Simulation und -Integration mit ROS, Gazebo, RViz, Docker, Ubuntu und NVIDIA Isaac Sim.
    \item Vertrautheit mit Reinforcement-Learning-Trainingsworkflows für Robotik- und simulationsbasierte autonome Systeme.
    \item Vertrautheit mit ONNX-Modellbereitstellung für portable KI-Modellaustausch- und Inferenz-Workflows.
    \item ANN-basierte ingenieurwissenschaftliche Analyseforschung mit MATLAB-orientierten Workflows.
\end{itemize}

%====================
\section*{Ausbildung}

\textbf{\color{PrimaryColor}Bachelor of Engineering in Mechatronik, Robotik und Automatisierungstechnik} \hfill 2016 -- 2020\\
UMPSA Universiti Malaysia Pahang Al-Sultan Abdullah

%====================
\section*{Veröffentlichung}

\textbf{\color{PrimaryColor}The Flexural Strength Prediction of Porous Cu-Sn-Ti Composites via Artificial Neural Networks}\\
The 6th International Conference on Robotics Intelligence and Applications, 2018\\
DOI: \href{https://link.springer.com/chapter/10.1007/978-981-13-8323-6_34}{10.1007/978-981-13-8323-6\_34}

%====================
\section*{Kurse \& Zertifizierungen}

\begin{itemize}[noitemsep, topsep=5pt]
    \item Niveau C1 in Englisch -- Busuu (Mai 2025)
    \item Niveau A1 in Deutscher Sprache -- Busuu (Nov 2024)
    \item Relationale Datenbanken mit SQL -- Manara (Nov 2024)
    \item Grundlagen der Datenanalyse mit Google Sheets -- Coursera (Okt 2024)
    \item Grundlagen von Datenbanken -- Manara (Okt 2024)
    \item Lua-Skripting: Komplettes Lua-Programmieren von Grund auf -- Udemy (Nov 2023)
    \item Basisschulung Datenschutz (EU) -- Mitarbeiterschule GmbH (Sep 2023)
    \item ROS2 als ROS1-Entwickler lernen und ROS-Projekte migrieren -- Udemy (Nov 2022)
    \item Level 2: Simulation \& Hardwareimplementierung für SLAM -- Udemy (Nov 2022)
    \item Einstieg in Azure DevOps Boards -- Coursera (Okt 2022)
    \item Kommandozeile in Linux -- Coursera (Sep 2022)
    \item Arduino Workshop -- Core Electronics (Sep 2022)
    \item Open-Source-Softwareentwicklungsmethoden -- The Linux Foundation (Okt 2021)
    \item Python-Skripting für DevOps-Spezialisierung -- LearnQuest (Sep 2021)
    \item Python-Skripting: Dateien, Vererbung und Datenbanken -- LearnQuest (Sep 2021)
    \item Webanwendungstechnologien und Django -- University of Michigan (Sep 2021)
    \item AWS Cloud Technical Essentials -- AWS (Aug 2021)
    \item DevOps und Build-Automatisierung mit Python -- LearnQuest (Aug 2021)
    \item Python-Skripting: Daten, Klassen und Sammlungen -- LearnQuest (Aug 2021)
    \item Konfigurationsmanagement und die Cloud -- Google (Aug 2021)
    \item Einführung in Git und GitHub -- Google (Jul 2021)
    \item Grundlagen des technischen Supports -- Google (Jul 2021)
    \item Fehlerbehebung und Debugging-Techniken -- Google (Jul 2021)
    \item Verwendung von Python zur Interaktion mit dem Betriebssystem -- Google (Jun 2021)
    \item Python-Crashkurs -- Google (Mai 2021)
    \item Malaysia Autonomous Vehicle Challenge (MAVeC'19) -- MARii (Nov 2018)
    \item The International Conference of Robotics Intelligence and Applications (RiTA) 2018 -- KIAST SG (2018)
    \item Boundless Innovation Week -- HPMM (Aug 2018)
    \item CAM-Workshop (Rechnerunterstützte Fertigung) -- The Institution of Engineers, Malaysia (Nov 2017)
    \item IEM MATLAB Workshop -- The Institution of Engineers, Malaysia (Okt 2017)
    \item Arduino für Einsteiger Workshop -- SPARKTECH Team at IIUM (Nov 2016)
    \item Online-Datenbankkurs: SCOPUS -- University Malaysia Pahang (Nov 2016)
\end{itemize}

%====================
\section*{Sprachen}

Arabisch: Muttersprache oder zweisprachige Kompetenz\\
Englisch: Berufliche Arbeitskompetenz\\
Türkisch: Grundkenntnisse\\
Deutsch: Grundkenntnisse\\
Malaiisch: Eingeschränkte Arbeitskompetenz

\end{document}